% Management in Global Markets – Business Plan (CaféOlé!)
% Format and colours match Global Management Project.pdf
% Compile: pdflatex business-plan-guide.tex

\documentclass[11pt,a4paper]{article}
\usepackage[utf8]{inputenc}
\usepackage[T1]{fontenc}
\usepackage[margin=2.5cm]{geometry}
\usepackage{enumitem}
\usepackage[table]{xcolor}
\usepackage{parskip}
\usepackage{booktabs}
\usepackage{fancyhdr}
\definecolor{eublue}{RGB}{0,51,153}
\definecolor{charcoal}{RGB}{51,51,51}
\definecolor{notegrey}{RGB}{90,90,90}
\definecolor{cafebrown}{RGB}{101,67,33}
\definecolor{sectionblue}{RGB}{0,51,153}

\pagestyle{fancy}
\fancyhf{}
\fancyfoot[C]{\small\textcolor{sectionblue}{\thepage}}
\renewcommand{\headrulewidth}{0pt}
\renewcommand{\footrulewidth}{0pt}

\begin{document}

% ========== FRONT PAGE (like Global Management Project.pdf) ==========
\thispagestyle{empty}
\begin{center}
\vspace*{1.5cm}
{\large Facultad de Geograf\'ia e Historia}\\[1.5em]
{\LARGE\bfseries\color{eublue} Global Management:}\\[0.3em]
{\LARGE\bfseries\color{eublue} Business Plan}\\[1em]
{\large Caf\'eOl\'e! Clubworking}\\[2em]
{\normalsize Chloe Borders, Ava Pelz, Jade Zhao}\\[0.5em]
{\small \copyright\ 2024 Caf\'eOl\'e. Produced by WSI. All rights reserved.}\\[1em]
{\normalsize Profesor Geoffrey Ditta}\\[0.3em]
{\normalsize Marzo 9, 2026}
\end{center}
\newpage
\pagestyle{fancy}

% ========== COPYRIGHT PAGE ==========
\section*{\color{sectionblue}Copyright and Legal Information}
\small
\textbf{Legal Identity:}\\
CAF\'E OL\'E SL (Sociedad Limitada)\\
Registered Address: Calle Fern\'andez de los R\'ios, 3, 28015 Madrid.

\textbf{Copyright Statement (Caf\'eOl\'e! Clubworking):}\\
All content, branding, designs, and text related to Caf\'eOl\'e! Clubworking business materials are the property of Caf\'e Ol\'e SL or their respective owners. ``Caf\'eol\'e \copyright\ 2024 Caf\'e Ol\'e SL. All rights reserved.''

\textbf{Website Copyright:}\\
https://cafeole-clubworking.com/\\
``\copyright\ 2024 Caf\'eOl\'e. Produced by WSI. All rights reserved.''

\textbf{Academic Project Copyright Statement:}\\
Copyright \copyright\ 2026 GBMTCafeRedesign. All rights reserved. This business plan and all supporting materials are original work created for academic purposes. No portion of this document may be reproduced, distributed, or transmitted in any form or by any means without prior written permission from the authors.

\textbf{Intellectual Property Notice:}\\
All logos, branding systems, mockups, marketing materials, and design concepts included in this business plan are original creations developed for this academic project. These concepts are not currently registered trademarks and are presented solely for educational evaluation purposes.
\normalsize
\newpage

% ========== TABLE OF CONTENTS (like PDF) ==========
\section*{\color{sectionblue}Table of Contents}
\begin{itemize}[leftmargin=*,nosep]
\item \textcolor{charcoal}{Introduction to Caf\'eOl\'e! Clubworking}
\item \textcolor{charcoal}{Executive Summary}
\item \textcolor{charcoal}{Company Description}
\item \textcolor{charcoal}{Business Model Canvas}
\item \textcolor{charcoal}{PESTLE Analysis}
\item \textcolor{charcoal}{SWOT Analysis}
\item \textcolor{charcoal}{Economic and Social Analysis}
\item \textcolor{charcoal}{Competitive Analysis}
\item \textcolor{charcoal}{Strategic Plan}
\item \textcolor{charcoal}{Contingency Plan}
\item \textcolor{charcoal}{Marketing Plan}
\end{itemize}
\newpage

% ========== INTRODUCTION (like PDF - prose format) ==========
\section*{\color{sectionblue}Introduction to Caf\'eOl\'e! Clubworking}

Caf\'eOl\'e! is an existing work cafe located in Chamber\'i, Madrid. It is an environment that provides features such as wifi, private workspaces, networking events, and affordable snacks and drinks for customers. Customers can either reserve a space online at the cafe or walk in. For five euros per hour, customers can work at the cafe and get unlimited coffee and snacks for as long as they pay. If they want to purchase a day pass, it is 25 euros, which includes Caf\'eOl\'e!'s perks for the entire day.

\textbf{Scope-objective:}\\
Provide a safe and reliable workspace for students, especially those studying abroad or on an Erasmus. Allow study abroad students to connect in a more professional and academic environment.

\textbf{Improvement:}\\
Currently, Caf\'eOl\'e! is geared towards professionals seeking a reliable work environment. We propose to make Caf\'eOl\'e! a more suitable place for students, especially those studying abroad or on an Erasmus program. This improvement will include special deals for students, like discounts on hourly work rates and monthly memberships (e.g. 30 euros per month). These discounts could be advertised at schools, and students could use their student cards at the cafe. Additionally, we propose to add events designated to connecting students to each other and other opportunities in Madrid. This will allow them to feel more at home in a new city, and it will allow them to connect with other students who are passionate about their academic and professional goals. It could also motivate students to do their schoolwork because they will be surrounded by other motivated peers.

\textbf{Competitive advantage:}\\
A competitive advantage of this improvement is a work cafe with students as their main focus. There are several work cafes in Madrid, but few, if any, of them cater to students, especially study abroad students. With this, a competitive advantage is customer loyalty. When students find community at a place like this, they are likely to come back. This increases both profit and reputation of the business. Partnering with local universities like Complutense could increase the number of students who regularly go to Caf\'eOl\'e!.

\vspace{0.5em}
\noindent{\color{sectionblue}\rule{\textwidth}{0.4pt}}

% ========== EXECUTIVE SUMMARY (approx. 1 page) ==========
\section*{\color{sectionblue}Executive Summary}

Caf\'eOl\'e! is an existing work cafe located in Chamber\'i, Madrid. It provides features such as wifi, private workspaces, networking events, and affordable snacks and drinks for customers. Customers can either reserve a space online at the cafe or walk in. For five euros per hour, customers can work at the cafe and get unlimited coffee and snacks for as long as they pay. If they want to purchase a day pass, it is 25 euros, which includes Caf\'eOl\'e!'s perks for the entire day. Right now, the cafe is geared primarily towards professionals seeking a reliable work environment.

We propose to shift the focus so that students can use it too, especially Erasmus and study abroad participants. They need somewhere to work that is not their flat, and they often want to meet others. Our improvement will include student discounts (20\% off the hourly rate), a monthly membership of 30 euros, university partnerships (Complutense), and community events such as language exchanges and study groups. Students can use their student cards at the cafe. The point is to give them a reliable spot and a bit of community. Study abroad can be lonely; when students find community at a place like this, they are likely to come back. This increases both profit and reputation of the business.

Madrid has plenty of cafes but not many that really work for students who need to sit with a laptop for hours. A lot of them have limited wifi, or they will ask you to put the laptop away when it gets busy, especially on weekends. Caf\'eOl\'e! is different. It is built for it. High-speed internet, quiet zones, communal seating. So there is a gap in the market and we are already set up for it. Few work cafes in Madrid cater specifically to students. Caf\'eOl\'e! can fill the gap.

PESTLE analysis indicates supportive conditions. Tourism is strong, with over 11 million visitors in 2025. Unemployment is at 9\%, the lowest since 2008. People want work-friendly cafes. Chamber\'i is right by University City. Spain's regulations (Sociedad Limitada, health certifications, labour laws) are clear. Social media and online reviews influence visits. Sustainable packaging aligns with younger generations. SWOT highlights our strengths: location, wifi, printers, community. Weaknesses: menu not on website, prices can add up. Opportunities: university partnerships, tourism, gentrification. Threats: La Bicicleta, Federal Cafe, seasonal tourism, rising costs.

Our competitive advantage is a work cafe with students as the main focus. La Bicicleta and Federal Cafe are established competitors, but the majority of cafes in Madrid have limited space, limited free wifi, and can turn away laptops at peak hours. Where they fall short, Caf\'eOl\'e! can step in. We have existing infrastructure: high-speed internet, dedicated quiet zones, communal seating. With further additions that help cater more toward a student population, Caf\'eOl\'e! has the potential to become the next best cafe that both students and professionals can thrive in.

Key risks include competitor response and cost pressures. We have contingency plans: if student uptake is low, we push university marketing and can refocus on professionals; if competitors copy us, we emphasise infrastructure and partnerships; if costs rise, we negotiate suppliers; if tourism drops, we target locals; if tech fails, we ensure backup internet. Success depends on execution. University partnerships take time to formalise. Student uptake must be validated. Overall, we conclude that the student-focused improvement is viable and will support customer loyalty, profitability, and reputation.

\vspace{0.5em}
\noindent{\color{sectionblue}\rule{\textwidth}{0.4pt}}

% ========== COMPANY DESCRIPTION ==========
\section*{\color{sectionblue}Company Description}

Caf\'eOl\'e! Clubworking is a Sociedad Limitada (SL) located in Chamber\'i, Madrid, near University City and corporate areas. Products and services include the work cafe model: 5 euros per hour or 25 euros day pass, unlimited coffee and snacks, wifi, private workspaces, printers, quiet zones. Student membership: 30 euros per month. Events: language exchange, study groups. Mission: provide reliable workspace and community for professionals and students. Values: quality coffee, reliable wifi, inclusive community. Value proposition: ``Feel at home in a new city''---safe, reliable workspace and community for study abroad students. Operations: Madrid; international clientele (Erasmus, study abroad); potential expansion.

\vspace{0.5em}
\noindent{\color{sectionblue}\rule{\textwidth}{0.4pt}}

% ========== BUSINESS MODEL CANVAS ==========
\section*{\color{sectionblue}Business Model Canvas}

\textbf{\textcolor{cafebrown}{Customer segments:}} Students (Erasmus, study abroad) and professionals seeking a reliable work environment.

\textbf{\textcolor{cafebrown}{Value proposition:}} Reliable wifi, community events (language exchange, study groups), unlimited coffee and snacks, quiet zones, printers. ``Feel at home in a new city.''

\textbf{\textcolor{cafebrown}{Channels:}} In-store, online reservations (cafeole-clubworking.com), social media (Instagram, TikTok).

\textbf{\textcolor{cafebrown}{Revenue streams:}} Hourly (5 euros), day pass (25 euros), student membership (30 euros/month), premium membership (up to 300 euros/month).

\textbf{\textcolor{cafebrown}{Key resources:}} Wifi, printers, space, quality coffee and snacks.

\textbf{\textcolor{cafebrown}{Key activities:}} Events (language exchange, study groups), university partnerships, community building.

\textbf{\textcolor{cafebrown}{Key partnerships:}} Complutense, local universities, potential collaborations with bakeries.

\textbf{\textcolor{cafebrown}{Cost structure:}} Rent (Chamber\'i 15--25\% of revenue), supplies (coffee, snacks 25--35\%), labour (minimum wage 1,134 euros full-time, 620 euros part-time), marketing (2--5\%).

\vspace{0.5em}
\noindent{\color{sectionblue}\rule{\textwidth}{0.4pt}}

% ========== PESTLE ANALYSIS (like PDF - prose) ==========
\section*{\color{sectionblue}PESTLE Analysis}

\textbf{\textcolor{cafebrown}{Political Factors:}} Following Spain's small business regulations, such as Sociedad Limitada, a limited liability company. City of Madrid policies and regulations for small businesses, such as Caf\'eOl\'e!. In Madrid, cafes must have health and hygiene certifications and a business operating license. The cafe must comply with Spain's strict labor laws, like minimum wage requirements and working hour regulations.

\textbf{\textcolor{cafebrown}{Economic Factors:}} Caf\'eOl\'e! must consider student budgets, as they typically have limited amounts of money. Many students are also unemployed, which would make the student discount helpful. They must also factor in the high cost of living in Madrid, which could decrease the number of customers. However, tourism in Madrid is high, which could increase profits. Tourism could bring travelers who need to work remotely.

\textbf{\textcolor{cafebrown}{Social Factors:}} Demand for reliable WiFi and work-friendly cafes in Madrid increases the popularity of places like Caf\'eOl\'e!. Many students are looking for good places to study near their university. Caf\'eOl\'e! is located close to the ``University City'' part of Madrid, which could increase customer numbers. Professionals in need of a change in environment from their work are driven to go to work cafes, increasing Caf\'eOl\'e!'s demand. It is prevalent among many people that cafes should have healthy food and drinks, which could influence demand.

\textbf{\textcolor{cafebrown}{Technological Factors:}} Online reviews through Yelp or Google could influence a customer on their decision to visit the cafe. Social media activity is another factor because an aesthetic Instagram or a popular TikTok about the cafe could sway a customer into visiting. Fast WiFi and reliable internet are key to retaining customers. Printers and contactless payment systems such as Apple Pay are also useful technological factors for the business.

\textbf{\textcolor{cafebrown}{Legal Factors:}} Caf\'eOl\'e! must adhere to national health regulations through the Agency for Food Safety and Nutrition (AESAN). HACCP (Hazard Analysis and Critical Control Points) is required for businesses to follow for food safety. They must also follow minimum wage requirements, which means that they must pay employees at least 1,134 euros a month if full-time (40 hours) and 620 euros a month for part-time (20 hours).

\textbf{\textcolor{cafebrown}{Environmental Factors:}} Brand image can be improved if they use sustainable packaging and practices. Younger generations generally pay attention to recycling and sustainable practices, which would help with the cafe's new focus on student customers. Additionally, Caf\'eOl\'e! could reduce environmental impact and long-term costs by investing in cleaner energy appliances.

\vspace{0.5em}
\noindent{\color{sectionblue}\rule{\textwidth}{0.4pt}}

% ========== SWOT ANALYSIS (like PDF - prose) ==========
\section*{\color{sectionblue}SWOT Analysis}

\textbf{\textcolor{cafebrown}{Strengths:}} Caf\'eOl\'e!'s ability to bring a reliable environment for a wide range of professional and student audiences. Quality coffee and pastries, all included in the hourly usage fee. Central location in Chamber\'i, close to corporate areas and universities. Fast WiFi and printers, variables that are not found in all coffee shops in Madrid. Community aspects like networking and language exchanges.

\textbf{\textcolor{cafebrown}{Weaknesses:}} Menu not available on the website, making it difficult for potential customers to determine their visit. Prices can get expensive for longer periods of time, and customers may not feel like they are getting their money's worth if they are not utilising the unlimited food and drink offer as much as they can. Dependence on in-store sales instead of delivery options.

\textbf{\textcolor{cafebrown}{Opportunities:}} Increased access to study abroad and Erasmus networks through partnerships with universities. Tourism in Madrid is experiencing increasing popularity, and the Chamber\'i neighbourhood is experiencing gentrification. Sustainable packaging and vegan options could increase customer interest. Collaborations with local and beloved bakeries to increase popularity.

\textbf{\textcolor{cafebrown}{Threats:}} Other coworking spaces in different locations. La Bicicleta in Malasa\~na and Federal Cafe in Chamber\'i are laptop and study-friendly competitors. Changes in consumer preferences in the coffee shop aesthetic or location. Fluctuations in tourism through different seasons. Fluctuating costs of resources like coffee beans, milk, and pastries from local bakeries.

\vspace{0.5em}
\noindent{\color{sectionblue}\rule{\textwidth}{0.4pt}}

% ========== ECONOMIC AND SOCIAL ANALYSIS (like PDF) ==========
\section*{\color{sectionblue}Economic and Social Analysis}

In 2025, Madrid had over 11 million tourists, meaning that tourism is a reliable source of customers for Caf\'eOl\'e!. Considering the number of tourists, many people want a reliable space where they can catch up on work, which is why Caf\'eOl\'e! is beneficial. Additionally, there is an increasing social trend of cafe culture, which is ideal for the cafe. Spain is facing inflation, as many countries are. This could increase the cost of food and cafe supplies, which could lead to price increases. Fortunately, Spain's unemployment rate was 9 percent in 2025, the lowest it has been since 2008. This allows for more customers, especially those who need a space to work.

\vspace{0.5em}
\noindent{\color{sectionblue}\rule{\textwidth}{0.4pt}}

% ========== COMPETITIVE ANALYSIS (like PDF - prose) ==========
\section*{\color{sectionblue}Competitive Analysis}

As discussed, Caf\'eOl\'e! caters itself as a ``clubworking'' and ``professional/coworking'' coffee destination with premium membership prices up to 300 euros per month with inclusion such as wifi, workstations, available coffee, printers, scanners, storage services, private rooms, and meeting rooms. However, 300 euros a month can be a steep price for students, as many do not work, isolating this cafe as a ``professional'' one. Moreover, Caf\'eOl\'e! faces stiff competition due to Madrid's competitive cafe scene and already-established cafes that are laptop and study-friendly, such as La Bicicleta in Malasa\~na and the Federal Cafe in Chamber\'i. For example, La Bicicleta offers moderate pricing for food and specialty coffees, laptop-friendly tables, access to wifi, and an urban space that makes this cafe highly competitive and extremely appealing to not only professionals, but the student population as well.

However, the majority of cafes in Madrid that cater to students and professionals looking to work with peers on technology tend to have limited space, limited free wifi, and individuals can be turned away or told to put their laptops away during peak hours of the day, especially during the weekends. These regulations and protocols can make it challenging for students and professionals looking for an affordable space to work and collaborate with one another at any time during the day, into the late afternoon. Where these cafes fall short, Caf\'eOl\'e! can step in to bridge that gap. Caf\'eOl\'e! has existing infrastructure such as high-speed internet, dedicated quiet zones, and communal seating. With further additions that help cater more toward a student population, Caf\'eOl\'e! has the potential to become the next best cafe that both students and professionals can thrive in.

\vspace{0.5em}
\noindent{\color{sectionblue}\rule{\textwidth}{0.4pt}}

% ========== STRATEGIC PLAN (like PDF) ==========
\section*{\color{sectionblue}Strategic Plan}

Objectives: (1) Make Caf\'eOl\'e! the go-to place for students and Erasmus in Chamber\'i. (2) Retain professionals. We are not replacing them, just adding. The strategy is differentiation. What makes us different: we focus on students, we do community events (language exchange, study groups, tutoring), we have wifi and printers that work, unlimited coffee and snacks. Other cafes in Madrid either limit laptop use or have unreliable wifi. We have the location (Chamber\'i, near universities, corporate, tourism), the infrastructure, the space.

Short-term (0--12 months): Roll out the student discount (20\% off hourly, 30 euros monthly membership). Put the menu on the website. Get in touch with Complutense for a partnership. Flyers, accept student cards for discount. Run 1 event per week at first (language exchange or study group). Scale to 2--3 per month if attendance justifies it.

Long-term (1--3 years): Be the recognised work cafe for study abroad students in Madrid. Once that is working, look at a second location or franchise. Key success factors: university partnerships, student retention, word-of-mouth. Maintain wifi, coffee, quiet zones.

\vspace{0.5em}
\noindent{\color{sectionblue}\rule{\textwidth}{0.4pt}}

% ========== CONTINGENCY PLAN (like PDF) ==========
\section*{\color{sectionblue}Contingency Plan}

\textbf{\textcolor{cafebrown}{Low student uptake:}} If after 6 months we have fewer than 15--20 regular student users, we push harder on university marketing. Complutense, flyers. Run a survey to find out what is not working. Adjust the offer. If it still does not work after 12 months, we refocus on professionals. We are not burning bridges.

\textbf{\textcolor{cafebrown}{Competitor response:}} La Bicicleta in Malasa\~na and Federal Cafe in Chamber\'i are the main competitors. If they start copying our student offer, we lean into what they cannot easily copy: the dedicated work setup, printers, quiet zones, university partnerships. We have the infrastructure.

\textbf{\textcolor{cafebrown}{Rising costs:}} Inflation hits food and supplies. We would talk to suppliers for better terms. Maybe a small price adjustment. We would try to keep the student rate reasonable.

\textbf{\textcolor{cafebrown}{Seasonal drop:}} Tourism peaks in summer, dips Jan--Feb. Erasmus arrivals Sept/Feb. When it is quiet, we focus more on locals and semester memberships for students so they are locked in.

\textbf{\textcolor{cafebrown}{Tech failure:}} If the wifi goes, we are in trouble. We would need backup internet and proper maintenance. Fast wifi and reliable internet are key for retention.

\vspace{0.5em}
\noindent{\color{sectionblue}\rule{\textwidth}{0.4pt}}

% ========== MARKETING PLAN (like PDF) ==========
\section*{\color{sectionblue}Marketing Plan}

\textbf{\textcolor{cafebrown}{Market research:}} Qualitative (interviews with students: what do they want? which events? willingness to pay?) and quantitative (survey: how many would use 30 euros per month? how often? demographics). Both required by the course.

\textbf{\textcolor{cafebrown}{Marketing Mix (4P's):}}
\begin{itemize}[leftmargin=*,nosep]
\item \textbf{Product:} Work cafe with wifi, printers, coffee, snacks, quiet zones, community events, student membership.
\item \textbf{Price:} 5 euros/hour, 25 euros day pass. Student: 30 euros/month. Premium: up to 300 euros/month.
\item \textbf{Place:} Chamber\'i, online reservations, walk-in.
\item \textbf{Promotion:} Instagram, TikTok, Yelp, Google, university flyers, Complutense partnership, word-of-mouth.
\end{itemize}

\textbf{\textcolor{cafebrown}{Calendar:}}
\begin{itemize}[leftmargin=*,nosep]
\item Month 1: Social media setup, menu on website, student discount announced.
\item Month 2: Flyers at Complutense, student card verification, first posts.
\item Month 3: First language exchange or study group.
\item Months 4--6: Add second event type if M3 works; university partnership talks.
\item Months 7--9: September new Erasmus cohort; push membership.
\item Months 10--12: Evaluate, adjust, plan Year 2.
\end{itemize}

\textbf{\textcolor{cafebrown}{Budget Year 1:}}

\begin{center}
\small
\begin{tabular}{@{}lr@{}}
\toprule
\rowcolor{sectionblue!15}
\textbf{\textcolor{sectionblue}{Line item}} & \textbf{\textcolor{sectionblue}{Euros}} \\
\midrule
Social media (Instagram, TikTok ads, content) & 600 \\
Flyers, print (Complutense, universities) & 150 \\
Event materials (language exchange, study groups) & 300 \\
Website (menu, updates) & 100 \\
Complutense partnership (in-kind or small fee) & 50 \\
\midrule
\textbf{Total marketing budget Year 1} & \textbf{1,200} \\
\bottomrule
\end{tabular}
\end{center}

\vspace{0.5em}
\noindent{\color{sectionblue}\rule{\textwidth}{0.4pt}}

\small\textcolor{notegrey}{\textbf{Sources:} Class notes weeks 1--7, Global Management Project.pdf, Business plan reunidas, Global MGMT Business Plan. No external sources.}

\end{document}
